% Content of Worksheet 01, shared by the problem sheet and the solution sheet.
% Solutions are wrapped in the `solution` environment and appear only when the
% driver defines \ipsshowsolutions.

\ipsmaketitle

This worksheet is ungraded practice, with the same shape and level as the final exam.
Plan up to 3 hours for a serious pass.
Work each problem through on paper before you open the solutions, and show your algebra:
the exam asks for the steps, not only the number.
Some parts state an intermediate result ("show that\,\ldots"); use it to check your work, and to continue if a step resists.
Carry at least four significant figures through the intermediate steps (the worked solutions print five) and round only at the end.

\begin{given}[Constants and data]
$G = 6.674 \times 10^{-11}\ \mathrm{m^3\,kg^{-1}\,s^{-2}}$ \quad
$1\ \mathrm{AU} = 1.496 \times 10^{11}\ \mathrm{m}$ \quad
$1\ \mathrm{yr} = 3.156 \times 10^{7}\ \mathrm{s}$ \quad
$1\ \mathrm{d} = 86\,400\ \mathrm{s}$ \quad
$\Msun = 1.989 \times 10^{30}\ \mathrm{kg}$ \quad
$\Mearth = 5.972 \times 10^{24}\ \mathrm{kg}$ \quad
$\Rearth = 6371\ \mathrm{km}$

\vskip6pt
\begin{tabular}{@{}lrrrr@{}}
\toprule
Body & $a$ & $P$ & Mass $(\Mearth)$ & $\rho\ (\mathrm{kg\,m^{-3}})$ \\
\midrule
Mercury & $0.387\ \mathrm{AU}$ & $0.241\ \mathrm{yr}$ & $0.055$ & $5427$ \\
Venus   & $0.723\ \mathrm{AU}$ & $0.615\ \mathrm{yr}$ & $0.815$ & $5243$ \\
Earth   & $1.000\ \mathrm{AU}$ & $1.000\ \mathrm{yr}$ & $1.000$ & $5514$ \\
Mars    & $1.524\ \mathrm{AU}$ & $1.881\ \mathrm{yr}$ & $0.107$ & $3934$ \\
Jupiter & $5.203\ \mathrm{AU}$ & $11.86\ \mathrm{yr}$ & $317.8$ & $1326$ \\
Saturn  & $9.537\ \mathrm{AU}$ & $29.46\ \mathrm{yr}$ & $95.16$ & $687$ \\
Uranus  & $19.19\ \mathrm{AU}$ & $84.01\ \mathrm{yr}$ & $14.54$ & $1271$ \\
Neptune & $30.07\ \mathrm{AU}$ & $164.8\ \mathrm{yr}$ & $17.15$ & $1638$ \\
\midrule
Moon      & $3.844 \times 10^{5}\ \mathrm{km}$ & $27.32\ \mathrm{d}$ & $0.0123$ & $3344$ \\
Io        & & $1.769\ \mathrm{d}$ & & \\
Europa    & & $3.551\ \mathrm{d}$ & & \\
Ganymede  & $1.0704 \times 10^{6}\ \mathrm{km}$ & $7.155\ \mathrm{d}$ & & \\
\bottomrule
\end{tabular}

\vskip4pt
{\footnotesize Saturn's mean radius is $58\,232\ \mathrm{km}$. Planetary and satellite values from
NASA/JPL Solar System Dynamics; Galilean satellite periods as quoted in the Lecture~2 notes.}
\end{given}


% ════════════════════════════════════════════════════════════
\problem{Weighing the Sun and Jupiter}

Newton's form of Kepler's third law turns a period and an orbital size into a mass.
Here you apply the same equation twice, to two different central bodies, and then ask how far it can be trusted.

\parti A planet of mass $m$ moves on a circular orbit of radius $r$ about a star of mass $M$,
with $m \ll M$. Starting from the balance between the gravitational force and the centripetal
force, derive
\[
  M = \frac{4\pi^2 r^3}{G P^2}.
\]

\begin{solution}
The gravitational force supplies the centripetal force:
\[
  \frac{G M m}{r^2} = \frac{m v^2}{r}.
\]
The orbital speed on a circular orbit is $v = 2\pi r / P$, so
\[
  \frac{G M m}{r^2} = \frac{m}{r}\left(\frac{2\pi r}{P}\right)^{2}
                    = \frac{m}{r}\,\frac{4\pi^2 r^2}{P^2}
                    = \frac{4\pi^2 m r}{P^2}.
\]
Divide both sides by $m$:
\[
  \frac{G M}{r^2} = \frac{4\pi^2 r}{P^2}.
\]
Multiply both sides by $r^2/G$:
\[
  M = \frac{4\pi^2 r^3}{G P^2}.
\]
\end{solution}

\parti Judge this claim in 2 to 3 sentences: \emph{"If Jupiter were twice as massive, its
orbital period at its present distance would be noticeably longer."}

\begin{solution}
False, on two counts. The planet's mass cancels in part~(a): it sets both the gravitational
pull and the inertia, so the period contains only the stellar mass and the orbital size, and
doubling Jupiter changes nothing at this order. In the exact two-body form,
$P^2 = 4\pi^2 a^3 / [G(M + m)]$, a heavier planet raises the total attracting mass, so the
period becomes slightly \emph{shorter}, not longer; for a second Jupiter mass the change is
$\tfrac{1}{2}\,\Mjup/\Msun \approx 5 \times 10^{-4}$, far below "noticeable".
\answer{False: to leading order the planet's mass cancels, and the exact correction is tiny
and has the opposite sign.}
\end{solution}

\parti Apply part~(a) twice. First compute the mass of the Sun from Jupiter's orbit
($a = 5.203\ \mathrm{AU}$, $P = 11.86\ \mathrm{yr}$; show first that
$r = 7.784 \times 10^{11}\ \mathrm{m}$ and $P = 3.743 \times 10^{8}\ \mathrm{s}$). Then compute
the mass of Jupiter from Ganymede's orbit ($a = 1.0704 \times 10^{6}\ \mathrm{km}$,
$P = 7.155\ \mathrm{d}$; show that $P^2 = 3.822 \times 10^{11}\ \mathrm{s^2}$), and express it in
Earth masses.

\begin{solution}
\textbf{The Sun from Jupiter.} Convert to SI units:
\[
  r = 5.203 \times 1.496 \times 10^{11}\ \mathrm{m} = 7.7837 \times 10^{11}\ \mathrm{m},
  \qquad
  P = 11.86 \times 3.156 \times 10^{7}\ \mathrm{s} = 3.7430 \times 10^{8}\ \mathrm{s}.
\]
Form the two powers:
\[
  r^3 = (7.7837 \times 10^{11})^3 = 4.7158 \times 10^{35}\ \mathrm{m^3},
  \qquad
  P^2 = (3.7430 \times 10^{8})^2 = 1.4010 \times 10^{17}\ \mathrm{s^2}.
\]
Numerator and denominator separately:
\[
  4\pi^2 r^3 = 39.478 \times 4.7158 \times 10^{35} = 1.8617 \times 10^{37}\ \mathrm{m^3},
\]
\[
  G P^2 = 6.674 \times 10^{-11} \times 1.4010 \times 10^{17}
        = 9.3503 \times 10^{6}\ \mathrm{m^3\,kg^{-1}}.
\]
Divide:
\[
  M = \frac{1.8617 \times 10^{37}}{9.3503 \times 10^{6}} = 1.991 \times 10^{30}\ \mathrm{kg}.
\]

\textbf{Jupiter from Ganymede.} The same equation, with Jupiter now the central body:
\[
  a = 1.0704 \times 10^{9}\ \mathrm{m},
  \qquad
  P = 7.155 \times 86\,400\ \mathrm{s} = 6.1819 \times 10^{5}\ \mathrm{s},
\]
\[
  a^3 = (1.0704 \times 10^{9})^3 = 1.2264 \times 10^{27}\ \mathrm{m^3},
  \qquad
  P^2 = (6.1819 \times 10^{5})^2 = 3.8216 \times 10^{11}\ \mathrm{s^2},
\]
\[
  M = \frac{4\pi^2 a^3}{G P^2}
    = \frac{39.478 \times 1.2264 \times 10^{27}}{6.674 \times 10^{-11} \times 3.8216 \times 10^{11}}
    = \frac{4.8416 \times 10^{28}}{25.505}
    = 1.898 \times 10^{27}\ \mathrm{kg}.
\]
In Earth masses:
\[
  \frac{1.898 \times 10^{27}}{5.972 \times 10^{24}} = 317.9.
\]
\answer{$\Msun \approx 1.991 \times 10^{30}\ \mathrm{kg}$ ($0.10\%$ above the accepted
$1.989 \times 10^{30}\ \mathrm{kg}$; part~(d) accounts for that excess) and
$\Mjup \approx 1.898 \times 10^{27}\ \mathrm{kg} = 317.9\,\Mearth$, matching the data table.}
\end{solution}

\parti Your solar mass in part~(c) exceeds the accepted value by $0.10\%$. Use the exact form
$P^2 = 4\pi^2 a^3 / [G(M + m)]$ to explain what that computation actually measured, and to
account for the excess. Would the same effect be visible in your Jupiter mass? The mass of
Ganymede is $1.482 \times 10^{23}\ \mathrm{kg}$.

\begin{solution}
The exact form returns the \emph{sum} of the two masses, so the first computation, which
neglected $m$, actually measured $\Msun + \Mjup$:
\[
  \Msun + \Mjup = \Msun\left(1 + \frac{\Mjup}{\Msun}\right),
  \qquad
  \frac{\Mjup}{\Msun} = \frac{1.898 \times 10^{27}}{1.991 \times 10^{30}}
                      = 9.533 \times 10^{-4},
\]
\[
  1.989 \times 10^{30} \times (1 + 9.533 \times 10^{-4}) = 1.9909 \times 10^{30}\ \mathrm{kg},
\]
which agrees with the $1.9911 \times 10^{30}\ \mathrm{kg}$ of part~(c) to $0.01\%$. The $0.10\%$
excess is therefore not an arithmetic slip and not a rounding artefact: it is Jupiter's own mass.

For Jupiter from Ganymede the same correction is
\[
  \frac{M_{\mathrm{Gan}}}{\Mjup} = \frac{1.482 \times 10^{23}}{1.898 \times 10^{27}}
                                 = 7.81 \times 10^{-5},
\]
that is $0.008\%$, a factor of $12$ smaller and below the four-significant-figure precision of
the input data.
\answer{The neglected orbiter mass shifts the derived central mass by the fraction $m/M$:
$0.095\%$ for Jupiter about the Sun, invisible $0.008\%$ for Ganymede about Jupiter.}
\end{solution}


% ════════════════════════════════════════════════════════════
\problem{The cheapest route to Mars}

Treat the orbits of Earth and Mars as circular and coplanar, with radii $1.000\ \mathrm{AU}$ and
$1.524\ \mathrm{AU}$. A spacecraft transfers between them on an ellipse whose perihelion touches
Earth's orbit and whose aphelion touches Mars's orbit. At each end the spacecraft fires its engine
briefly (a \emph{burn}), changing its speed by an increment $\Delta v$; a burn along the direction
of motion is \emph{prograde}. Work entirely in the Sun's gravitational field and ignore the
gravity of Earth and Mars themselves.

\parti Two spacecraft of equal mass are on orbits with the same semi-major axis $a$ but different
eccentricities. Judge each claim, with justification: (i) \emph{"they have the same total
energy"}; (ii) \emph{"they pass through the distance $r = a$ at the same speed"}.

\begin{solution}
Both claims are true. (i) The total energy of a Keplerian orbit is $E = -GMm/2a$, which depends
on $a$ alone; eccentricity only changes how the fixed total is divided between kinetic and
potential energy along the orbit. (ii) Put $r = a$ into the vis-viva equation:
\[
  v^2 = GM\left(\frac{2}{a} - \frac{1}{a}\right) = \frac{GM}{a},
  \qquad\text{so}\qquad
  v = \sqrt{\frac{GM}{a}},
\]
independent of $e$. On the circular orbit this speed holds everywhere; on the eccentric orbit it
holds at the two points where the body is at distance $a$ from the focus, the ends of the minor
axis.
\answer{Both true; both follow from $E$ depending only on $a$.}
\end{solution}

\parti Show that the transfer ellipse has $a_t = 1.262\ \mathrm{AU}$. Explain why the flight
time is half the period of that ellipse, and compute it in days.

\begin{solution}
The transfer ellipse has perihelion $r_p = 1.000\ \mathrm{AU}$ and aphelion
$r_a = 1.524\ \mathrm{AU}$. Since $r_p = a(1 - e)$ and $r_a = a(1 + e)$, adding the two gives
\[
  r_p + r_a = a(1 - e) + a(1 + e) = 2a ,
\]
so
\[
  a_t = \frac{r_p + r_a}{2} = \frac{1.000 + 1.524}{2} = 1.262\ \mathrm{AU}
      = 1.88795 \times 10^{11}\ \mathrm{m}.
\]
Departure is at perihelion and arrival is at aphelion. These are the two ends of the major axis, so
they are separated by half a revolution and the flight takes half a period.

With $GM = 6.674 \times 10^{-11} \times 1.989 \times 10^{30}
      = 1.32746 \times 10^{20}\ \mathrm{m^3\,s^{-2}}$:
\[
  a_t^3 = (1.88795 \times 10^{11})^3 = 6.7293 \times 10^{33}\ \mathrm{m^3},
  \qquad
  \frac{a_t^3}{GM} = \frac{6.7293 \times 10^{33}}{1.32746 \times 10^{20}}
                   = 5.0693 \times 10^{13}\ \mathrm{s^2},
\]
\[
  P_t = 2\pi \sqrt{5.0693 \times 10^{13}}
      = 2\pi \times 7.1199 \times 10^{6}
      = 4.4736 \times 10^{7}\ \mathrm{s}.
\]
Half of that is
\[
  t = \tfrac{1}{2} P_t = 2.2368 \times 10^{7}\ \mathrm{s}
    = \frac{2.2368 \times 10^{7}}{86\,400} = 258.9\ \mathrm{d}.
\]
\answer{$a_t = 1.262\ \mathrm{AU}$; flight time $259\ \mathrm{d}$, about $8.5$ months.}

A faster route to the same number: dividing Kepler's third law for the transfer ellipse by the
same law for Earth's orbit ($a = 1\ \mathrm{AU}$, $P = 1\ \mathrm{yr}$) cancels every constant
and leaves $P^2 = a^3$ in units of AU and years, so $P_t = 1.262^{3/2} = 1.418\ \mathrm{yr}$ and
$t = 0.709\ \mathrm{yr}$.
\end{solution}

\parti Use the vis-viva equation to show that the spacecraft moves at
$v_p = 32.7\ \mathrm{km\,s^{-1}}$ at the transfer perihelion and $v_a = 21.5\ \mathrm{km\,s^{-1}}$
at the transfer aphelion. Then compute the circular orbital speeds of Earth and Mars, and from the
four speeds the two velocity increments $\Delta v_1$ (departure), $\Delta v_2$ (arrival), and
their sum.

\begin{solution}
At perihelion, $r = 1.000\ \mathrm{AU} = 1.496 \times 10^{11}\ \mathrm{m}$ and
$a = a_t = 1.88795 \times 10^{11}\ \mathrm{m}$:
\[
  \frac{2}{r} = 1.3369 \times 10^{-11}\ \mathrm{m^{-1}},
  \qquad
  \frac{1}{a_t} = 5.2968 \times 10^{-12}\ \mathrm{m^{-1}},
  \qquad
  \frac{2}{r} - \frac{1}{a_t} = 8.0722 \times 10^{-12}\ \mathrm{m^{-1}},
\]
\[
  v_p^2 = GM\left(\frac{2}{r} - \frac{1}{a_t}\right)
        = 1.32746 \times 10^{20} \times 8.0722 \times 10^{-12}
        = 1.0716 \times 10^{9}\ \mathrm{m^2\,s^{-2}},
\]
\[
  v_p = 3.2735 \times 10^{4}\ \mathrm{m\,s^{-1}} = 32.73\ \mathrm{km\,s^{-1}}.
\]
At aphelion, $r = 1.524\ \mathrm{AU} = 2.2799 \times 10^{11}\ \mathrm{m}$:
\[
  \frac{2}{r} = 8.7723 \times 10^{-12}\ \mathrm{m^{-1}},
  \qquad
  \frac{2}{r} - \frac{1}{a_t} = 3.4755 \times 10^{-12}\ \mathrm{m^{-1}},
\]
\[
  v_a^2 = 1.32746 \times 10^{20} \times 3.4755 \times 10^{-12}
        = 4.6136 \times 10^{8}\ \mathrm{m^2\,s^{-2}},
\]
\[
  v_a = 2.1479 \times 10^{4}\ \mathrm{m\,s^{-1}} = 21.48\ \mathrm{km\,s^{-1}}.
\]
The circular speeds are the $r = a$ case of the vis-viva equation:
\[
  v_\oplus = \sqrt{\frac{GM}{r_p}}
           = \sqrt{\frac{1.32746 \times 10^{20}}{1.496 \times 10^{11}}}
           = \sqrt{8.8734 \times 10^{8}}
           = 29.79\ \mathrm{km\,s^{-1}},
\]
\[
  v_{\mathrm{Mars}} = \sqrt{\frac{1.32746 \times 10^{20}}{2.2799 \times 10^{11}}}
                    = \sqrt{5.8224 \times 10^{8}}
                    = 24.13\ \mathrm{km\,s^{-1}}.
\]
The increments:
\[
  \Delta v_1 = v_p - v_\oplus = 32.73 - 29.79 = 2.95\ \mathrm{km\,s^{-1}},
\]
\[
  \Delta v_2 = v_{\mathrm{Mars}} - v_a = 24.13 - 21.48 = 2.65\ \mathrm{km\,s^{-1}},
\]
\[
  \Delta v_{\mathrm{total}} = 2.95 + 2.65 = 5.60\ \mathrm{km\,s^{-1}}.
\]
\answer{$\Delta v_1 = 2.95$, $\Delta v_2 = 2.65$, total $\Delta v = 5.60\ \mathrm{km\,s^{-1}}$.}
\end{solution}

\parti Sketch the spacecraft's speed $v(r)$ along the transfer, from departure at
$1.000\ \mathrm{AU}$ to arrival at $1.524\ \mathrm{AU}$. Mark the four speeds of part~(c): $v_p$,
$v_a$, and the circular speeds of Earth and Mars. Use the sketch to explain why \emph{both} burns
point in the direction of motion.

\begin{solution}
The vis-viva speed falls monotonically from perihelion to aphelion:

\begin{center}
\begin{tikzpicture}[line cap=round]
  % axes
  \draw[->, ipsText] (0,0) -- (6.4,0) node[right=1pt, font=\footnotesize] {$r$ (AU)};
  \draw[->, ipsText] (0,0) -- (0,4.4) node[above=2pt, font=\footnotesize] {$v$ (km\,s$^{-1}$)};
  % x ticks at 1.000, 1.262, 1.524 AU  [x = (r-0.95)*9]
  \foreach \x/\lab in {0.45/1.000, 2.81/1.262, 5.17/1.524}{
    \draw[ipsText] (\x,0) -- (\x,-0.08) node[below, font=\scriptsize] {\lab};}
  % y ticks at 20, 25, 30 km/s  [y = (v-20)*0.30]
  \foreach \y/\lab in {0.0/20, 1.5/25, 3.0/30}{
    \draw[ipsText] (0,\y) -- (-0.08,\y) node[left, font=\scriptsize] {\lab};}
  % transfer curve through computed points
  \draw[ipsAccent, thick] plot[smooth] coordinates
    {(0.45,3.82) (0.90,3.43) (1.35,3.05) (2.25,2.36) (3.15,1.72) (4.05,1.13) (5.17,0.44)};
  % endpoint markers
  \fill[ipsAccent] (0.45,3.82) circle (0.055) node[above right=1pt, font=\scriptsize, color=ipsText] {$v_p = 32.7$};
  \fill[ipsAccent] (5.17,0.44) circle (0.055) node[below left=1pt, font=\scriptsize, color=ipsText] {$v_a = 21.5$};
  % circular speeds
  \draw[ipsOrange, fill=white, thick] (0.45,2.94) circle (0.055) node[right=3pt, font=\scriptsize, color=ipsText] {$v_\oplus = 29.8$};
  \draw[ipsOrange, fill=white, thick] (5.17,1.24) circle (0.055) node[above left=1pt, font=\scriptsize, color=ipsText] {$v_{\mathrm{Mars}} = 24.1$};
  % burn arrows
  \draw[->, ipsOrange, thick] (0.45,2.94) -- (0.45,3.76);
  \node[right, font=\scriptsize, color=ipsOrange] at (0.55,3.36) {$\Delta v_1$};
  \draw[->, ipsOrange, thick] (5.17,0.50) -- (5.17,1.18);
  \node[left, font=\scriptsize, color=ipsOrange] at (5.07,0.84) {$\Delta v_2$};
\end{tikzpicture}
\end{center}

At departure the transfer demands $32.7\ \mathrm{km\,s^{-1}}$ while Earth's orbit provides
$29.8$: the spacecraft must \emph{gain} speed, so the first burn is prograde. At arrival the
spacecraft crawls at its aphelion speed of $21.5\ \mathrm{km\,s^{-1}}$ while Mars passes at
$24.1$: the spacecraft is the slower body and must again \emph{gain} speed to match, so the
second burn is prograde too.
\answer{The transfer curve starts above Earth's circular speed and ends below Mars's; closing
each gap means adding speed along the direction of motion, so both burns are prograde.}
\end{solution}


% ════════════════════════════════════════════════════════════
\problem{The Laplace resonance of the Galilean moons}

Lecture~2 describes the orbital periods of Io, Europa and Ganymede as standing in a $1:2:4$ ratio.
This problem asks how exact that statement is, and what the resonance geometry enforces.

\parti Compute the two consecutive period ratios and the Ganymede-to-Io ratio, and give the
percentage by which each misses its integer value.

\begin{solution}
\[
  \frac{P_{\mathrm{Eur}}}{P_{\mathrm{Io}}} = \frac{3.551}{1.769} = 2.00735,
  \qquad
  \frac{2.00735 - 2}{2} = +0.37\% ,
\]
\[
  \frac{P_{\mathrm{Gan}}}{P_{\mathrm{Eur}}} = \frac{7.155}{3.551} = 2.01493,
  \qquad
  \frac{2.01493 - 2}{2} = +0.75\% ,
\]
\[
  \frac{P_{\mathrm{Gan}}}{P_{\mathrm{Io}}} = \frac{7.155}{1.769} = 4.04466,
  \qquad
  \frac{4.04466 - 4}{4} = +1.12\% .
\]
\answer{The ratios exceed $2$, $2$ and $4$ by $0.37\%$, $0.75\%$ and $1.12\%$.}

These offsets are real, not measurement error: the periods are known to far better than a part in
$10^{4}$. Part~(b) shows what the resonance does hold exactly: a combination of all three mean
motions, not any pairwise ratio.
\end{solution}

\parti The mean motion is $n = 2\pi / P$. Evaluate the combination
\[
  n_{\mathrm{Io}} - 3\,n_{\mathrm{Eur}} + 2\,n_{\mathrm{Gan}}
\]
and express the result as a fraction of $n_{\mathrm{Io}}$. Keep 7 decimal digits through the
arithmetic; as a checkpoint, $3\,n_{\mathrm{Eur}} = 0.8448324 \times 2\pi\ \mathrm{d^{-1}}$.
Contrast the outcome with part~(a).

\begin{solution}
The factor $2\pi$ is common to all three terms, so it can be carried outside and the arithmetic done
with $1/P$ in units of $\mathrm{d^{-1}}$:
\[
  \frac{1}{P_{\mathrm{Io}}}  = \frac{1}{1.769} = 0.5652911\ \mathrm{d^{-1}},
\]
\[
  \frac{3}{P_{\mathrm{Eur}}} = \frac{3}{3.551} = 3 \times 0.2816108 = 0.8448324\ \mathrm{d^{-1}},
\]
\[
  \frac{2}{P_{\mathrm{Gan}}} = \frac{2}{7.155} = 2 \times 0.1397624 = 0.2795248\ \mathrm{d^{-1}}.
\]
Combine, keeping every digit because this is a difference of nearly equal numbers:
\[
  0.5652911 - 0.8448324 = -0.2795413\ \mathrm{d^{-1}},
\]
\[
  -0.2795413 + 0.2795248 = -0.0000165\ \mathrm{d^{-1}},
\]
so, restoring the $2\pi$,
\[
  n_{\mathrm{Io}} - 3 n_{\mathrm{Eur}} + 2 n_{\mathrm{Gan}}
    = -1.04 \times 10^{-4}\ \mathrm{rad\,d^{-1}},
  \qquad
  \frac{n_{\mathrm{Io}} - 3 n_{\mathrm{Eur}} + 2 n_{\mathrm{Gan}}}{n_{\mathrm{Io}}}
    = -2.9 \times 10^{-5}.
\]
\answer{The combination vanishes to three parts in $10^{5}$, although no individual ratio in
part~(a) came closer than three parts in $10^{3}$.}

This is the Laplace relation. Each pairwise ratio misses its integer by a few tenths of a percent,
yet the particular three-body combination is zero to a hundred times better precision. The exact
constraint the system obeys is the one on all three moons together, not the one on any pair. With
the NASA/JPL sidereal periods quoted to seven digits ($1.769138$, $3.551181$, $7.154553\ \mathrm{d}$) the residual
falls further, to about $10^{-7}$ of $n_{\mathrm{Io}}$.
\end{solution}

\parti The figure shows the three orbits at the moment of a Europa--Ganymede conjunction; the
resonance locks the phases such that Io then stands on the \emph{opposite} side of Jupiter.

\begin{center}
\begin{tikzpicture}[line cap=round]
  \draw[gray!55] (0,0) circle (0.62);
  \draw[gray!55] (0,0) circle (1.00);
  \draw[gray!55] (0,0) circle (1.52);
  \fill[ipsPrimary] (0,0) circle (0.10);
  \node[below, font=\scriptsize, color=ipsText] at (0,-0.14) {Jupiter};
  \fill[ipsOrange] (180:0.62) circle (0.065);
  \node[left=2pt, font=\scriptsize, color=ipsText] at (180:0.62) {Io};
  \fill[ipsAccent] (0:1.00) circle (0.065);
  \node[above=2pt, font=\scriptsize, color=ipsText] at (0:1.00) {Europa};
  \fill[ipsPrimary!75] (0:1.52) circle (0.065);
  \node[right=2pt, font=\scriptsize, color=ipsText] at (0:1.52) {Ganymede};
\end{tikzpicture}
\end{center}

Idealise the periods as exactly $1:2:4$. Sketch the positions of the three moons after $1$, $2$,
and $4$ Io periods. From your sketch: (i) where does Io stand at every Europa--Ganymede
conjunction, and (ii) can a triple conjunction occur?

\begin{solution}
Per Io period, Io advances $360^\circ$, Europa $180^\circ$, and Ganymede $90^\circ$. Starting from
the printed configuration (Io at $180^\circ$, Europa and Ganymede at $0^\circ$):

\begin{center}
\begin{tabular}{@{}lccc@{}}
\toprule
Time (Io periods) & Io & Europa & Ganymede \\
\midrule
$0$ & $180^\circ$ & $0^\circ$   & $0^\circ$   \\
$1$ & $180^\circ$ & $180^\circ$ & $90^\circ$  \\
$2$ & $180^\circ$ & $0^\circ$   & $180^\circ$ \\
$4$ & $180^\circ$ & $0^\circ$   & $0^\circ$   \\
\bottomrule
\end{tabular}
\end{center}

\begin{center}
\begin{tikzpicture}[line cap=round, scale=0.62]
  \foreach \dx/\t/\io/\eu/\ga in {0/0/180/0/0, 4.2/1/180/180/90, 8.4/2/180/0/180, 12.6/4/180/0/0}{
    \begin{scope}[shift={(\dx,0)}]
      \draw[gray!55] (0,0) circle (0.62);
      \draw[gray!55] (0,0) circle (1.00);
      \draw[gray!55] (0,0) circle (1.52);
      \fill[ipsPrimary] (0,0) circle (0.09);
      \fill[ipsOrange] (\io:0.62) circle (0.09);
      \fill[ipsAccent] (\eu:1.00) circle (0.09);
      \fill[ipsPrimary!75] (\ga:1.52) circle (0.09);
      \node[below, font=\scriptsize, color=ipsText] at (0,-1.75) {$t = \t\,P_{\mathrm{Io}}$};
    \end{scope}}
  \node[right, font=\scriptsize, color=ipsText, align=left] at (14.5,0.6)
    {\textcolor{ipsOrange}{$\bullet$} Io\\[-1pt]
     \textcolor{ipsAccent}{$\bullet$} Europa\\[-1pt]
     \textcolor{ipsPrimary!75}{$\bullet$} Ganymede};
\end{tikzpicture}
\end{center}

Because Io always advances by full turns per Io period, it returns to $180^\circ$ at every integer
time; Europa alternates between $0^\circ$ and $180^\circ$; Ganymede steps through
$90^\circ$ per Io period. The pattern repeats after $4$ Io periods, one Ganymede period.

(i) Europa--Ganymede conjunctions occur at $t = 0, 4, 8, \ldots$, always at $0^\circ$, and Io then
stands at $180^\circ$: at every Europa--Ganymede conjunction, Io is on the opposite side of
Jupiter. (ii) Each pairwise conjunction finds the third moon elsewhere: Io--Europa conjunctions
(odd $t$, at $180^\circ$) find Ganymede at $90^\circ$ or $270^\circ$; the Io--Ganymede conjunction
($t = 2$, at $180^\circ$) finds Europa at $0^\circ$.
\answer{A triple conjunction never occurs. The resonance phase-protects the system: the moons
never line up all on one side, which caps how coherently their mutual perturbations can add.}
\end{solution}

\parti The resonance holds Io's eccentricity at $e \approx 0.004$ rather than letting tides
circularise the orbit. Explain why a non-zero eccentricity is needed for tidal heating to continue
in a moon that is already tidally locked, and name the observable consequence.

\begin{solution}
A tidally locked moon on a perfectly circular orbit keeps the same face toward the planet at a
constant distance. Its tidal bulge is then fixed in the body frame: it neither rises and falls nor
moves across the surface, so no work is done deforming the interior and no heat is released.

An eccentric orbit breaks this in two ways at once. The distance to the planet varies over each
orbit, so the height of the bulge varies with it. And the moon spins at a constant rate while its
orbital angular velocity is fastest at periapsis and slowest at apoapsis, so the direction to the
planet librates back and forth in longitude and drags the bulge with it. Both effects flex the
interior once per orbit, and internal friction converts that mechanical work into heat.

Left alone, tidal dissipation would damp the eccentricity to zero and switch itself off. The
Laplace resonance prevents that by continually re-forcing $e$.
\answer{Observable consequence: Io is the most volcanically active body in the solar system.}
\end{solution}


% ════════════════════════════════════════════════════════════
\problem{Tides and the Roche limit}

\parti A moon of radius $R_s$ orbits at distance $d$ from a planet of mass $M_p$, with
$R_s \ll d$. Show that the difference between the planet's gravitational acceleration at the moon's
near surface and at its centre is
\[
  \Delta a \approx \frac{2 G M_p R_s}{d^3}.
\]
Use $(1 - x)^{-2} \approx 1 + 2x$ for $x \ll 1$.

\begin{solution}
The near surface sits at distance $d - R_s$ from the planet, the centre at $d$:
\[
  a_{\mathrm{near}} = \frac{G M_p}{(d - R_s)^2},
  \qquad
  a_{\mathrm{centre}} = \frac{G M_p}{d^2}.
\]
Factor $d$ out of the first denominator:
\[
  a_{\mathrm{near}} = \frac{G M_p}{d^2}\left(1 - \frac{R_s}{d}\right)^{-2}
                    \approx \frac{G M_p}{d^2}\left(1 + \frac{2 R_s}{d}\right).
\]
Subtract:
\[
  \Delta a = a_{\mathrm{near}} - a_{\mathrm{centre}}
           = \frac{G M_p}{d^2}\left(1 + \frac{2 R_s}{d}\right) - \frac{G M_p}{d^2}
           = \frac{G M_p}{d^2} \cdot \frac{2 R_s}{d}
           = \frac{2 G M_p R_s}{d^3}.
\]
The direct pull falls as $d^{-2}$ but the tidal difference falls as $d^{-3}$, so their ratio is
$\Delta a / a = 2 R_s / d$: tides are a small correction far away and dominate close in.
\end{solution}

\parti The fluid Roche limit is
\[
  d_R \approx 2.46\, R_p \left(\frac{\rho_p}{\rho_s}\right)^{1/3}.
\]
Interrogate this formula before using it: (i) why does the \emph{satellite's} radius $R_s$ not
appear? (ii) what does it predict as $\rho_s \to \infty$, and does that make sense? (iii) what
does it predict for $\rho_s = \rho_p$, and what does that say about how close \emph{any}
self-gravitating satellite can orbit?

\begin{solution}
(i) Both sides of the comparison scale the same way with satellite size: the tidal acceleration
across the satellite, $2 G M_p R_s / d^3$, grows linearly with $R_s$, and so does the satellite's
own surface gravity, $G m_s / R_s^2 = \tfrac{4}{3}\pi G \rho_s R_s$. The radius cancels in the
ratio and only the densities survive. Big and small satellites of the same density are torn apart
at the same distance.

(ii) As $\rho_s \to \infty$, $d_R \to 0$: an arbitrarily dense satellite can orbit arbitrarily
close, because its self-gravity grows without bound while the tidal stress at fixed $d$ does not.
The trend is physically sensible.

(iii) For $\rho_s = \rho_p$ the density factor is unity and $d_R = 2.46\,R_p$: even a satellite
as dense as its planet is destroyed within $2.46$ planetary radii. The limit falls inside the
planet, $d_R < R_p$, only for $\rho_s > 2.46^3\,\rho_p \approx 15\,\rho_p$, far denser than any
satellite material, so for every real satellite density the forbidden zone extends well beyond
the planet's surface.
\answer{$R_s$ cancels because tidal stress and self-gravity both scale linearly with it; only the
density ratio matters, and even equal density leaves a forbidden zone of $2.46\,R_p$.}
\end{solution}

\parti Evaluate $d_R$ for Saturn ($R_p = 58\,232\ \mathrm{km}$,
$\rho_p = 687\ \mathrm{kg\,m^{-3}}$) and a satellite of solid water ice
($\rho_s = 1000\ \mathrm{kg\,m^{-3}}$); as a checkpoint, $(687/1000)^{1/3} = 0.8824$. Saturn's
main rings span $67\,000$ to $137\,000\ \mathrm{km}$ from the planet's centre, with the F ring
near $140\,000\ \mathrm{km}$: compare. Ring particles are in fact porous aggregates with
$\rho_s \approx 500\ \mathrm{kg\,m^{-3}}$; \emph{without} recomputing from scratch, state the
factor by which the limit moves and what that does to your comparison.

\begin{solution}
Solid ice:
\[
  d_R = 2.46 \times 58\,232 \times 0.8824 = 2.46 \times 51\,384
      = 1.264 \times 10^{5}\ \mathrm{km} = 2.17\,R_p .
\]
This sits at $126\,400\ \mathrm{km}$: the C, B and inner A rings lie inside it, but the outer A
ring ($126\,400$ to $137\,000\ \mathrm{km}$) and the F ring lie outside, which the solid-ice
limit fails to explain.

Halving the density multiplies $d_R$ by $(1000/500)^{1/3} = 2^{1/3} = 1.26$, because
$d_R \propto \rho_s^{-1/3}$:
\[
  d_R = 1.26 \times 126\,400 = 1.59 \times 10^{5}\ \mathrm{km} = 2.73\,R_p ,
\]
and the \emph{entire} ring system, F ring included, now lies inside the limit.
\answer{$d_R = 126\,400\ \mathrm{km}$ ($2.17\,R_p$) for solid ice; porosity moves it out by
$2^{1/3} = 1.26$ to $159\,000\ \mathrm{km}$ ($2.73\,R_p$), enclosing the whole ring system.}
\end{solution}

\parti For the Earth--Moon pair the same formula gives $d_R \approx 18\,500\ \mathrm{km}
= 2.91\,\Rearth$; the Moon orbits $21$ times further out. The International Space Station orbits
$420\ \mathrm{km}$ above Earth's surface, far \emph{inside} that limit. Judge this claim in 2 to
3 sentences: \emph{"the ISS survives inside the Roche limit, so the Roche formula must be
wrong."}

\begin{solution}
False. The Roche limit compares the tidal stress with a body's \emph{self-gravity}, so it applies
only to objects held together by gravity alone: fluid bodies and loose aggregates bound only by their own self-gravity. The ISS is held
together by the material strength of its structure, which at its mass exceeds its self-gravity by
many orders of magnitude, so it sits outside the formula's domain rather than refuting it. The
same escape applies to solid rock, which is why small rigid moons survive well inside the fluid
limit.
\answer{False: the formula presupposes self-gravity as the only glue; the ISS is bolted, not
gravitationally bound.}
\end{solution}


% ════════════════════════════════════════════════════════════
\problem{From planetesimals to planets, and the mass of the solar nebula}

\parti Gravitational focusing raises the collision cross-section of a body of radius $R$ to
\[
  \sigma_{\mathrm{eff}} = \pi R^2 \left(1 + \frac{v_{\mathrm{esc}}^2}{v_\infty^2}\right).
\]
Two planetesimals of density $3000\ \mathrm{kg\,m^{-3}}$ have radii $50\ \mathrm{km}$ and
$500\ \mathrm{km}$. Show that the smaller body has $v_{\mathrm{esc}} = 64.8\ \mathrm{m\,s^{-1}}$,
and explain why $v_{\mathrm{esc}} \propto R$ at fixed density. Then compute both focusing factors
for a swarm with random velocity $v_\infty = 10\ \mathrm{m\,s^{-1}}$, and compare the ratio of the
two effective cross-sections with the ratio of the geometric ones.

\begin{solution}
For the $50\ \mathrm{km}$ body, $R = 5.0 \times 10^{4}\ \mathrm{m}$:
\[
  M = \tfrac{4}{3}\pi R^3 \rho
    = \tfrac{4}{3}\pi (5.0 \times 10^{4})^3 \times 3000
    = 1.571 \times 10^{18}\ \mathrm{kg},
\]
\[
  v_{\mathrm{esc}} = \sqrt{\frac{2 G M}{R}}
    = \sqrt{\frac{2 \times 6.674 \times 10^{-11} \times 1.571 \times 10^{18}}{5.0 \times 10^{4}}}
    = \sqrt{4194}
    = 64.8\ \mathrm{m\,s^{-1}} .
\]
At fixed density $M \propto R^3$, so
$v_{\mathrm{esc}} = \sqrt{2GM/R} \propto \sqrt{R^3/R} = R$: the $500\ \mathrm{km}$ body has ten
times the escape speed, $v_{\mathrm{esc}} = 648\ \mathrm{m\,s^{-1}}$.

Focusing factors at $v_\infty = 10\ \mathrm{m\,s^{-1}}$, using
$v_{\mathrm{esc}}^2 = 4194\ \mathrm{m^2\,s^{-2}}$ from above and, since
$v_{\mathrm{esc}} \propto R$, $100$ times that value for the larger body:
\[
  1 + \frac{4194}{10^2} = 1 + 41.9 = 42.9
  \qquad\text{and}\qquad
  1 + \frac{4.194 \times 10^{5}}{10^2} = 1 + 4194 = 4195 .
\]
The cross-sections are in the ratio
\[
  \frac{\sigma_{\mathrm{eff}}(500)}{\sigma_{\mathrm{eff}}(50)}
    = \left(\frac{500}{50}\right)^2 \times \frac{4195}{42.9}
    = 100 \times 97.8
    = 9.8 \times 10^{3},
\]
against a geometric ratio of only $(500/50)^2 = 100$.
\answer{Focusing gives the larger body an extra factor of $98$ on top of its geometric advantage.}
\end{solution}

\parti Read the limits of the focusing formula: what does $\sigma_{\mathrm{eff}}$ reduce to for
$v_\infty \gg v_{\mathrm{esc}}$, and how does it scale with $R$ for
$v_\infty \ll v_{\mathrm{esc}}$? Confirm the first limit by evaluating both focusing factors at
$v_\infty = 500\ \mathrm{m\,s^{-1}}$. The growing bodies themselves stir the swarm toward their
own escape speed: name the two growth regimes these limits correspond to, and explain why the
first regime shuts itself down.

\begin{solution}
For $v_\infty \gg v_{\mathrm{esc}}$ the bracket tends to $1$ and
$\sigma_{\mathrm{eff}} \to \pi R^2$: the purely geometric cross-section, growing only as $R^2$.
For $v_\infty \ll v_{\mathrm{esc}}$ the second term dominates,
\[
  \sigma_{\mathrm{eff}} \approx \pi R^2 \frac{v_{\mathrm{esc}}^2}{v_\infty^2}
    \propto R^2 \times R^2 = R^4,
\]
using $v_{\mathrm{esc}} \propto R$ from part~(a): the biggest body gains a steeply
super-geometric advantage.

At $v_\infty = 500\ \mathrm{m\,s^{-1}}$, with the same $v_{\mathrm{esc}}^2$ values:
\[
  1 + \frac{4194}{500^2} = 1.02,
  \qquad
  1 + \frac{4.194 \times 10^{5}}{500^2} = 2.68,
\]
so the large body's advantage over pure $R^2$ scaling collapses from the factor $98$ of part~(a)
to $2.6$.

The cold-swarm limit is \textbf{runaway growth}: the largest body eats fastest, pulls away from
the swarm, and the mass spread widens. The hot-swarm limit is \textbf{oligarchic growth}: a few
similar-sized bodies grow slowly and in step. The transition is self-inflicted: the winners of
runaway growth gravitationally stir the planetesimals, raising $v_\infty$ toward their own
$v_{\mathrm{esc}}$, which switches their focusing advantage off.
\answer{Geometric $\pi R^2$ for a hot swarm, $\propto R^4$ for a cold one; runaway growth creates
the stirring that ends it, handing over to oligarchic growth.}
\end{solution}

\parti Estimate the minimum mass of the solar nebula. Assume the four terrestrial planets are made
entirely of solids, each of the four giant planets needed a solid core of $15\,\Mearth$, and
solids make up $1\%$ of the disk by mass. As a checkpoint, the terrestrial planets contribute
$1.98\,\Mearth$ of solids. Give the answer in Earth masses and as a fraction of $\Msun$.

\begin{solution}
Solids locked into the terrestrial planets:
\[
  0.055 + 0.815 + 1.000 + 0.107 = 1.977\ \Mearth .
\]
Solids locked into the four giant-planet cores:
\[
  4 \times 15 = 60\ \Mearth .
\]
Total solid inventory:
\[
  M_{\mathrm{solid}} = 1.977 + 60 = 61.98\ \Mearth .
\]
If solids are $1\%$ of the disk by mass, the disk must have held
\[
  M_{\mathrm{disk}} = \frac{61.98}{0.01} = 6198\ \Mearth
    = 6198 \times 5.972 \times 10^{24} = 3.701 \times 10^{28}\ \mathrm{kg},
\]
\[
  \frac{M_{\mathrm{disk}}}{\Msun} = \frac{3.701 \times 10^{28}}{1.989 \times 10^{30}} = 0.0186 .
\]
\answer{About $6.2 \times 10^{3}\,\Mearth$, or $1.9\%$ of a solar mass.}

This is a lower bound twice over. Giant-planet heavy-element inventories are probably larger than
a bare $15\,\Mearth$ core, and no disk converts $100\%$ of its solids into planets, so the true
requirement is higher.
\end{solution}

\parti By the Class~II stage of disk evolution, a typical disk holds about $0.1\%$ of the stellar
mass. Judge this claim, with numbers: \emph{"a Class~II disk still holds enough material to build
the solar system's planets."} What does your verdict imply for \emph{when} the planets' solid
material had to be locked into large bodies?

\begin{solution}
A Class~II disk around the Sun holds
\[
  0.001 \times \frac{1.989 \times 10^{30}}{5.972 \times 10^{24}} = 333\ \Mearth
\]
of gas and dust in total. The claim fails on both budgets. The four giant planets alone sum to
\[
  317.8 + 95.16 + 14.54 + 17.15 = 444.7\ \Mearth ,
\]
more than the entire disk. And at $1\%$ dust, the disk carries only $3.33\ \Mearth$ of solids
against the $61.98\ \Mearth$ required by part~(c), short by a factor of
\[
  \frac{61.98}{3.33} = 18.6 .
\]
\answer{False: a Class~II disk holds neither the gas ($333 < 445\,\Mearth$) nor the solids
($3.3$ vs $62\,\Mearth$) that the planets demand.}

The implication is about timing. The solid inventory, and the giant-planet cores with their gas
envelopes, had to be assembled \emph{earlier}, while the disk was still young and held a few
percent of a solar mass, as in part~(c). This is consistent with giant-planet cores completing
their growth within the first few Myr, before the gas disperses.
\end{solution}
